\documentclass[11pt]{article}
\usepackage[utf8]{inputenc}
\usepackage{amsmath, amssymb, amsthm}
\usepackage{hyperref}
\usepackage{booktabs}
\usepackage{graphicx}

\title{Efficient Threshold ML-DSA\\ (Integration Notes for luxfi/threshold)}
\author{Sof\'ia Celi \and Rafael del Pino \and Thomas Espitau \and Guilhem Niot \and Thomas Prest}
\date{USENIX Security 2026 — integrated into luxfi/threshold as \texttt{protocols/mldsa}}

\begin{document}
\maketitle

\begin{abstract}
This document records the reference for the Threshold ML-DSA paper
(Celi, del Pino, Espitau, Niot, Prest, USENIX Security 2026) as integrated
into \texttt{github.com/luxfi/threshold}. The implementation lives in
\texttt{protocols/mldsa/} and provides the first practical threshold
signature scheme fully compatible with NIST FIPS 204 (ML-DSA).

Output signatures are byte-compatible with standard ML-DSA, enabling drop-in
replacement of classical threshold ECDSA/Schnorr wallets with a post-quantum
scheme that keeps the standardized verification path.
\end{abstract}

\section{Summary of Construction}

The scheme tailors the template of Finally! (del Pino \& Niot, PKC 2025) to
ML-DSA's specific constraints:

\begin{itemize}
  \item \textbf{Short Replicated Secret Sharing (RSS).} For $(T,N)$, sample
        $\binom{N}{N-T+1}$ ML-DSA secrets $s_I \leftarrow \chi_s$ and distribute
        each $s_I$ to every party in $I$. The final public key is
        $\mathrm{vk} = \lfloor A\sum_I s_I \rceil$, matching the ML-DSA public
        key shape.
  \item \textbf{Unbalanced hyperball rejection.} Replace ML-DSA's uniform
        rejection with per-party hyperball rejection. The first $\ell$
        coordinates (no hint check) accept a wider ball than the last $k$
        coordinates (hint check), parameterized by $\nu > 1$.
  \item \textbf{Per-party + combination rejection.} Two rejection stages:
        (i) each party checks its local partial response has bounded norm;
        (ii) the combiner checks the aggregated $z^{(1)}$ satisfies
        $\|z^{(1)}\|_\infty < \gamma_1 - \beta$ and the hint is within the
        ML-DSA bound.
  \item \textbf{$K$ parallel instances.} To mitigate the probabilistic abort
        at $T$ parties, the protocol runs $K$ parallel sessions and outputs
        the first successful one. This amortizes the cost at the expense of
        bandwidth.
  \item \textbf{Optimized share reconstruction.} The RSS partition minimizes
        $\max_i |m_i|$ (number of secrets per party per session) via a
        max-flow on the bipartite \texttt{users × secrets} graph.
        For $N \leq 6$ the optimal partitions are hardcoded.
\end{itemize}

\section{Rounds and Communication}

Per-attempt signing requires three rounds (plus a one-shot commit phase):
\begin{enumerate}
  \item Each party samples $r_i \leftarrow \chi_r$, computes
        $w_i = A \cdot r_i$, publishes $\mathrm{cmt}_i = H_\mathrm{cmt}(\mathrm{vk}, i, w_i)$.
  \item Reveal $w_i$. All parties derive
        $\tilde c = H(\mu \| \mathrm{HighBits}(w, 2\gamma_2))$,
        $c = \mathrm{SampleInBall}(\tilde c)$.
  \item Compute local response $z_i = c \cdot s^{\mathrm{part}}_i + r_i$,
        apply hyperball rejection, publish $z_i^{(1)}$.
\end{enumerate}

Per-party communication at security level ML-DSA-44 (Table~3 of the paper):

\begin{center}
\begin{tabular}{rrrrrr}
\toprule
$N \setminus T$ & 2 & 3 & 4 & 5 & 6 \\
\midrule
2 & 10.5\,kB & & & & \\
3 & 15.8\,kB & 21.0\,kB & & & \\
4 & 15.8\,kB & 36.8\,kB & 42.0\,kB & & \\
5 & 15.8\,kB & 73.5\,kB & 157.4\,kB & 84.0\,kB & \\
6 & 21.0\,kB & 99.8\,kB & 388.4\,kB & 524.8\,kB & 194.2\,kB \\
\bottomrule
\end{tabular}
\end{center}

\section{Security}

Unforgeability reduces tightly to (i) the unforgeability of standard ML-DSA
and (ii) $\mathrm{MLWE}_{q,k,\ell,\chi}$ for
$\chi \in \{\chi_s, \chi_r, \chi_z\}$ (Theorem~3.2 of the paper). Static
dishonest-majority security is proven in the ROM; for $N \leq 6$ adaptive
security follows via complexity leveraging with at most 5 bits of loss.

\section{Integration Scope}

The luxfi/threshold package \texttt{protocols/mldsa} implements the scheme
on top of \texttt{cloudflare/circl/sign/mldsa} and reuses
\texttt{luxfi/lattice} primitives already present in this tree:

\begin{itemize}
  \item \texttt{protocols/mldsa/params.go} — parameter sets for
        ML-DSA-44/65/87, ports Tables 3, 10, 11.
  \item \texttt{protocols/mldsa/rss.go} — replicated secret sharing,
        with hardcoded optimal partitions for $N \leq 6$ per Appendix~B.
  \item \texttt{protocols/mldsa/hrej.go} — imbalanced hyperball rejection
        (Fig.~4 of the paper).
  \item \texttt{protocols/mldsa/keygen.go} — centralized keygen (Fig.~5)
        and DKG (Appendix~D).
  \item \texttt{protocols/mldsa/sign.go} — three-round signing protocol
        (Fig.~6).
  \item \texttt{protocols/mldsa/combine.go} — combine + verify (Fig.~7),
        producing standard ML-DSA signatures.
  \item \texttt{protocols/mldsa/a\_posteriori.go} — a~posteriori key
        sharing (\S3.3, Appendix~E) for migrating an existing ML-DSA key.
\end{itemize}

\section{References}

Celi, S., del Pino, R., Espitau, T., Niot, G., Prest, T.
\emph{Efficient Threshold ML-DSA}. USENIX Security Symposium, 2026.
Artifact: \texttt{doi.org/10.5281/zenodo.17963721}.

Reference implementation builds on CIRCL
(\texttt{github.com/cloudflare/circl}) and Lattigo
(\texttt{github.com/tuneinsight/lattigo}, mirrored as \texttt{github.com/luxfi/lattice}).

\end{document}
